%% bare_jrnl.tex
%% V1.4b
%% 2015/08/26
%% by Michael Shell
%% see http://www.michaelshell.org/
%% for current contact information.
%%
%% This is a skeleton file demonstrating the use of IEEEtran.cls
%% (requires IEEEtran.cls version 1.8b or later) with an IEEE
%% journal paper.
%%
%% Support sites:
%% http://www.michaelshell.org/tex/ieeetran/
%% http://www.ctan.org/pkg/ieeetran
%% and
%% http://www.ieee.org/
%%
%% ---------------------------------------------------------------
%% Modified by Abstrack (TFG).
%% Plantilla de referencia: IEEE Journal LaTeX Template oficial
%% (IEEEtran.cls, distribuido vía template-selector.ieee.org),
%% con el front-matter ajustado a las Author Guidelines de
%% IEEE Internet of Things Journal (ieee-iotj.org/guidelines-for-authors).
%% Abstract e introducción se insertan automáticamente.
%% ---------------------------------------------------------------

%%*************************************************************************
%% Legal Notice:
%% This code is offered as-is without any warranty either expressed or
%% implied; without even the implied warranty of MERCHANTABILITY or
%% FITNESS FOR A PARTICULAR PURPOSE!
%% User assumes all risk.
%% In no event shall the IEEE or any contributor to this code be liable for
%% any damages or losses, including, but not limited to, incidental,
%% consequential, or any other damages, resulting from the use or misuse
%% of any information contained here.
%%
%% All comments are the opinions of their respective authors and are not
%% necessarily endorsed by the IEEE.
%%
%% This work is distributed under the LaTeX Project Public License (LPPL)
%% ( http://www.latex-project.org/ ) version 1.3, and may be freely used,
%% distributed and modified. A copy of the LPPL, version 1.3, is included
%% in the base LaTeX documentation of all distributions of LaTeX released
%% 2003/12/01 or later.
%% Retain all contribution notices and credits.
%% ** Modified files should be clearly indicated as such, including  **
%% ** renaming them and changing author support contact information. **
%%*************************************************************************


\documentclass[journal]{IEEEtran}

\hyphenation{op-tical net-works semi-conduc-tor}


\begin{document}

\title{[ARTICLE TITLE]}

\author{[FIRST AUTHOR],~\IEEEmembership{Member,~IEEE,}
        and~[CO-AUTHOR]%
\thanks{[AFFILIATION DATA].}%
\thanks{Manuscript received [DATE].}}

\markboth{IEEE INTERNET OF THINGS JOURNAL,~Vol.~[VOL], No.~[NUM], [MONTH]~[YEAR]}%
{[AUTHORS]: [SHORT TITLE]}

\maketitle

% IEEE IoT-J: el abstract debe tener entre 150 y 250 palabras y no debe
% incluir citas bibliográficas ni ecuaciones.
\begin{abstract}
% [Pendiente de generacion]
\end{abstract}

\begin{IEEEkeywords}
[KEYWORD 1], [KEYWORD 2], [KEYWORD 3]
\end{IEEEkeywords}

\IEEEpeerreviewmaketitle


\section{Introduction}
\label{sec:intro}

La traducción automática se sitúa en el cruce de la inteligencia artificial y el procesamiento del lenguaje natural, y tiene como objetivo fundamental la conversión de texto entre diferentes idiomas. Este campo es de vital importancia, ya que las traducciones automáticas posibilitan la comunicación intercultural y el acceso a la información a escala global, abordando barreras lingüísticas que, de otro modo, limitarían la colaboración y el intercambio de conocimientos. En aplicaciones prácticas, tales como en empresas multinacionales, plataformas de contenido digital y servicios de atención al cliente, la traducción automática se vuelve imprescindible.

A lo largo de los años, numerosas técnicas han sido desarrolladas para abordar este desafío. Desde los primeros sistemas basados en reglas hasta los enfoques estadísticos que dieron paso a la era de la traducción automática neural, el progreso ha sido notable. Sin embargo, los modelos más recientes, incluyendo redes neuronales profundas y mecanismos de atención, han demostrado ser particularmente eficaces al permitir una mejor comprensión del contexto y la semántica de las oraciones complejas. A pesar de estos avances, todavía existen limitaciones significativas en la capacidad de estos modelos para manejar la variabilidad y la riqueza de las lenguas naturales en todas sus formas.

A pesar de los avances en la traducción automática mediante arquitecturas como redes neuronales profundas y modelos de atención, persisten desafíos significativos relacionados con el aprendizaje de dependencias a largo plazo y la adaptabilidad a oraciones de longitud variable. Por ejemplo, Bahdanau et al. (2014)~\cite{bahdanau2014neuralmt} y Sutskever et al. (2014)~\cite{sutskever2014seq2seq} han abordado estos problemas utilizando métodos de atención; sin embargo, aún se enfrentan a limitaciones inherentes a la codificación de oraciones en un solo vector de longitud fija. Esta limitación puede resultar en la incapacidad de captar adecuadamente el contexto necesario para traducir oraciones complejas y largas, lo que indica la necesidad de enfoques que eviten la saturación de información en un único vector.

Además, se observa que los modelos existentes suelen requerir un alto costo computacional y no siempre logran generalizar de manera óptima en diversos contextos lingüísticos, lo que resalta la falta de progreso en términos de eficiencia y efectividad en comparación con los sistemas de traducción basados en frases que, aunque más limitados en términos de flexibilidad, mantienen un rendimiento competitivo. Estos aspectos crean un contexto propicio para la investigación de nuevas metodologías que superen las barreras actuales y mejoren la calidad y velocidad de la traducción automática.

En respuesta a las limitaciones identificadas en trabajos previos sobre traducción automática, la propuesta central de este artículo es la introducción del modelo Transformador, el cual se basa exclusivamente en mecanismos de atención. Este enfoque permite manejar de manera más efectiva las dependencias de largo alcance, superando la limitación crítica de la codificación de oraciones completas en un vector de longitud fija. Al hacerlo, el Transformador evita los problemas que surgen de la saturación de información en un único vector, permitiendo capturar mejor el contexto necesario para traducciones precisas de oraciones complejas y largas.

Además, la arquitectura del Transformador facilita una mayor paralelización durante el entrenamiento, lo que no solo mejora la eficiencia de la formación del modelo, sino que también reduce significativamente el tiempo requerido para su entrenamiento. Este enfoque, en contraste con los modelos existentes que a menudo enfrentan dificultades en la generalización a diferentes contextos lingüísticos y son computacionalmente intensivos, promete no solo mejorar la calidad de las traducciones, sino también ofrecer una solución más escalable para aplicaciones prácticas en el campo de la traducción automática.

Las principales contribuciones del trabajo son las siguientes:
- Introducción del modelo Transformador, que se basa exclusivamente en mecanismos de atención y se describe en la sección 3 del documento.
- Superación de los resultados anteriores en tareas de traducción del conjunto de datos WMT 2014 con puntuaciones BLEU mejoradas, lo que se detalla en la sección 6.
- Validación de la robustez del modelo frente a la longitud de las oraciones, garantizando así su eficacia en una variedad de contextos y escenarios de traducción.

La validación del trabajo se lleva a cabo a través de una serie de experimentos en tareas de traducción automática, donde se utilizan conjuntos de datos estándar para evaluar el desempeño del modelo Transformador. Los experimentos permiten medir la calidad de las traducciones generadas mediante la puntuación BLEU, proporcionando una métrica objetiva que respalda las contribuciones del estudio. Se implementan análisis detallados que contrastan el modelo propuesto con arquitecturas anteriores, así como la observación de su robustez en la traducción de oraciones de diversas longitudes, lo que añade una capa adicional de evidencia sobre la efectividad del enfoque. Estas metodologías experimentales aseguran que las mejoras en el rendimiento no solo sean significativas, sino también consistentes a través de diferentes contextos de uso en el campo de la traducción automática.

El resto del documento se organiza en secciones que abordan de manera sistemática la introducción del contexto y el problema, los antecedentes sobre el estado del arte, la descripción detallada del modelo y sus componentes, seguido de la presentación de resultados experimentales, discusiones sobre los hallazgos y las conclusiones del estudio. Esta estructura permite una comprensión clara del desarrollo del trabajo, desde la motivación inicial hasta los resultados y su relevancia en el campo de la traducción automática.


\section{[SECTION 2 TITLE]}
\label{sec:section2}

[Section 2 content.]


\section{[SECTION 3 TITLE]}
\label{sec:section3}

[Section 3 content.]


\section{[SECTION 4 TITLE]}
\label{sec:section4}

[Section 4 content.]


\section{Conclusion}
\label{sec:conclusion}

[Conclusions.]


\appendices
\section{Appendix}
[Appendix content.]


\section*{Acknowledgment}

[Acknowledgments.]


\bibliographystyle{IEEEtran}
\bibliography{attention_is_all_you_need}


\begin{IEEEbiographynophoto}{[FIRST AUTHOR]}
[Brief author biography.]
\end{IEEEbiographynophoto}


\end{document}
